% ============================================================
%  Chapter 2 — Literature Review
% ============================================================
\section{Literature Review}
\label{ch:literature}

This chapter reviews the academic and regulatory literature relevant to the
thesis. \Cref{sec:lit:reg} covers the \IMO{} \CII{} regulatory framework and
the EU emissions trading system. \Cref{sec:lit:speed} reviews the speed
optimisation literature. \Cref{sec:lit:constrained} addresses
carbon-constrained operational planning models. \Cref{sec:lit:shortsea}
discusses the characteristics of short-sea and \RoRo{} shipping relevant
to this study. \Cref{sec:lit:gap} summarises the research gap.

% ── 2.1 Regulatory framework ──────────────────────────────────
\subsection{IMO Carbon Intensity Regulation}
\label{sec:lit:reg}

\subsubsection{Overview of the CII Framework}
\label{sec:lit:cii_overview}

The \CII{} framework was introduced through \IMO{} resolution
MEPC.336(76) \parencite{mepc336} and entered into force on 1~January 2023.
The framework requires each vessel to compute an \textit{attained \CII{}}
as the ratio of annual CO$_2$ emissions to a transport work proxy, compare
it against a year-specific \textit{required \CII{}}, and receive a letter
rating from A (best) to E (worst).  The required \CII{} for each year is
obtained by applying a reduction factor to a vessel-type-specific reference
line; the mandatory factors increase annually --- 5\,\% in 2023, 7\,\% in
2024, 9\,\% in 2025, and 11\,\% in 2026 --- thereby tightening the standard
each year.  Vessels receiving D or E ratings must submit corrective action
plans to their flag state; persistent D or E performance can trigger port
state control inspections and commercial consequences.

\subsubsection{CII Formula for RoRo Vessels}
\label{sec:lit:cii_roro}

The capacity measure used in the \CII{} denominator differs by ship type.
For most vessel types --- including bulk carriers and container ships ---
transport work is expressed as deadweight tonnage times nautical miles
sailed. For \RoRo{} vessels, however, the \IMO{} specifies Gross Tonnage
(\GT{}) as the capacity measure \parencite{mepc336, mepc339}. The attained
\CII{} for a \RoRo{} vessel is therefore:

\begin{equation}
  \ciiatt
  = \frac{\sum_k m^{\mathrm{CO}_2}_k}
         {f_i \cdot f_m \cdot \GT \cdot D}
  \label{eq:cii_roro}
\end{equation}

\noindent where $m^{\mathrm{CO}_2}_k$ is the mass of CO$_2$ emitted by
fuel type $k$, $\GT{}$ is the vessel's gross tonnage, $D$ is the total
distance sailed in nautical miles, and $f_i$, $f_m$ are correction factors
for ice class and capacity, respectively \parencite{mepc336}. The critical
implication of \cref{eq:cii_roro} is that cargo volume does not appear in
the denominator: \textit{a \RoRo{} vessel's \CII{} rating is independent
of how much cargo it carries.} Speed and distance are therefore the only
operational levers available to the operator.

\subsubsection{Rating Boundaries}
\label{sec:lit:boundaries}

\IMO{} resolution MEPC.339(76) specifies vessel-type-specific rating
boundary multipliers $d_1 < d_2 < d_3 < d_4$, such that a vessel receives
rating A if $\ciiatt < d_1 \cdot \ciireq$, rating B if
$d_1 \leq \ciiatt < d_2 \cdot \ciireq$, and so forth
\parencite{mepc339}. For \RoRo{} vessels, the boundaries are:
$d_1 = 0.76$, $d_2 = 0.89$, $d_3 = 1.08$, $d_4 = 1.27$.

It is important to note that the RoRo-specific boundary multipliers
($d_1=0.76$, $d_2=0.89$, $d_3=1.08$, $d_4=1.27$) differ significantly from
the generic values occasionally cited in the literature
(0.82\,/\,0.94\,/\,1.06\,/\,1.18).  Using the generic rather than the
vessel-type-specific boundaries would incorrectly assign letter ratings: for
the vessel in this study, the difference between correct and incorrect
boundaries shifts the Rating~A/B threshold from 10.94 to 11.78~\gGTnm{} ---
a non-trivial difference at the margins.  This thesis uses the RoRo-specific
multipliers throughout, consistent with MEPC.339(76) \parencite{mepc339}.

\subsubsection{Post-2026 Regulatory Trajectory}
\label{sec:lit:post2026}

Reduction factors beyond 2026 have not yet been formally adopted.
\IMO{} resolution MEPC.377(80) establishes a fleet-wide ambition of at
least 40\% reduction in GHG intensity by 2030 relative to 2008
\parencite{mepc377}. The precise annual per-vessel reduction factors for
2027--2030 are subject to review at MEPC~83 and~84 (2025--2026).
For this thesis, a 31\% per-vessel reduction by 2030 is assumed as a
representative intermediate scenario, acknowledging that the actual
outcome remains uncertain.

\subsubsection{EU Emissions Trading System}
\label{sec:lit:ets}

The European Union extended its Emissions Trading System (\ETS{}) to the
maritime sector from 1~January 2024, applying to all vessels of 5{,}000~GT
and above calling at EU ports \parencite{eu_ets_shipping}.  The phase-in
schedule is graduated: 40\,\% of verified emissions were covered by \ETS{}
in 2024, rising to 70\,\% in 2025 and 100\,\% from 2026 onwards.  Operators
must surrender EU Allowances (EUAs) corresponding to their covered emissions;
the EUA price is determined by the EU carbon market.  Historical EUA prices
have ranged from EUR~25 to EUR~100 per tonne of CO$_2$, with market
projections suggesting prices in the EUR~80--120/t range by 2030
\parencite{eu_ets_shipping}.

The EU \ETS{} and the \IMO{} \CII{} framework create a dual regulatory
pressure on Baltic Sea operators.  The \CII{} framework constrains the carbon
intensity of operations through a rating mechanism, while the \ETS{} attaches
a direct financial cost to every tonne of CO$_2$ emitted.  Both instruments
incentivise speed reduction, but they operate through different channels: \CII{}
through the risk of regulatory non-compliance, and \ETS{} through the voyage
cost function.  An operator facing high EUA prices and a tight \CII{} target
simultaneously faces compounding cost pressures that cannot be separated in
operational planning.  This interaction motivates the inclusion of an \ETS{}
cost term in the objective function of the present model (Scenarios S4 and S4b).

% ── 2.2 Speed optimisation ────────────────────────────────────
\subsection{Speed Optimisation in Maritime Shipping}
\label{sec:lit:speed}

Speed optimisation --- often called \textit{slow steaming} when speeds
are reduced below design --- is one of the most extensively studied
operational levers in maritime logistics
\parencite{psaraftis2013speed, ronen1982ship}.
The fundamental relationship is that fuel consumption scales
approximately as the cube of sailing speed:

\begin{equation}
  f(v) \approx \alpha v^3
  \label{eq:cubic_fuel}
\end{equation}

\noindent making speed reduction highly effective for emissions abatement.
\textcite{psaraftis2013speed} provide a comprehensive survey of speed
optimisation models and identify the trade-off between fuel savings and
transit time as the central economic tension.

The interaction between speed and regulatory constraints introduces important
complications.  In deep-sea tramp shipping, where schedule flexibility is
high, the cubic relationship \cref{eq:cubic_fuel} provides strong incentives
for slow steaming under both fuel cost and carbon intensity pressure
\parencite{PsaraftisKontovas2013SpeedModels}.  However, in fixed-schedule
liner services --- particularly short-sea services with tight port windows ---
the operator's ability to reduce speed is constrained by timetable commitments.
\textcite{ronen1982ship} identified this tension early, noting that the
optimal speed in a fixed-schedule context may be dictated by the schedule
rather than by cost minimisation.

Recent literature has begun to address \CII{}-aware speed optimisation
explicitly.  \textcite{cheng2025cii} demonstrate that treating the \CII{}
constraint as an ex-post metric rather than an operational constraint leads
to systematically suboptimal decisions.  \textcite{GolsefidiSharifiGhader2025WeatherSpeedFuel}
extend the speed optimisation problem to include weather uncertainty and
time-varying fuel quality, highlighting that the cubic approximation
\cref{eq:cubic_fuel} may require calibration for specific vessel types.
In the short-sea \RoRo{} context analysed in this thesis, the published
schedule places binding lower bounds on arc-level speeds, fundamentally
altering the feasible region of the speed optimisation problem.

% ── 2.3 Carbon-constrained planning ───────────────────────────
\subsection{Carbon-Constrained Operational Planning}
\label{sec:lit:constrained}

\textcite{cheng2025cii} develop a \MILP{} formulation for joint routing,
speed selection, and cargo acceptance under \CII{} constraints,
demonstrating that integrating the regulatory constraint into the
optimisation problem yields substantially different decisions compared
to treating \CII{} as an ex-post performance metric. Their formulation
serves as the methodological starting point for this thesis.  A key feature
of \textcite{cheng2025cii} is the linearisation of the cubic fuel function
via a speed grid, which makes the problem tractable for commercial \MILP{}
solvers while preserving the essential non-linearity of fuel cost.  Their
case study focuses on a tramp shipping context with flexible routing, where
the binding constraint is the \CII{} target rather than a schedule.

\textcite{GuWangIris2025AllianceUncertainFuel} address the related problem of
alliance planning under fuel uncertainty, demonstrating that the interaction
between fuel price volatility and carbon pricing materially affects optimal
fleet deployment.  While their focus is on container shipping, the insight that
\ETS{} costs and fuel costs co-determine optimal speeds is directly relevant
to the present analysis.

\textcite{GolsefidiSharifiGhader2025WeatherSpeedFuel} consider weather-dependent
speed and fuel consumption, showing that ignoring environmental variability can
lead to significant ex-post \CII{} violations.  Their findings motivate careful
calibration of the fuel model to the specific operating environment, as is done
in this thesis using EU \MRV{} data for the Baltic Sea route.

Across this body of work, the \CII{} constraint is typically modelled as a hard
annual constraint on the ratio of total CO$_2$ to total transport work.  The
present thesis adopts this approach but adds the dimension of schedule
feasibility, showing that the published timetable --- rather than the \CII{}
limit --- can be the operative binding constraint.

% ── 2.4 Short-sea and RoRo shipping ───────────────────────────
\subsection{Short-Sea and RoRo Shipping}
\label{sec:lit:shortsea}

Short-sea shipping --- typically defined as maritime transport not crossing
oceanic boundaries --- plays a critical role in European freight logistics,
with the Baltic Sea among the most active short-sea regions.  Baltic Sea
\RoRo{} services are characterised by high port call frequency (often daily
or bi-daily arrivals at multiple ports), relatively short individual legs
(200--1{,}200~nm), and fixed published timetables agreed with cargo shippers.
These structural features distinguish Baltic Sea operations from the deep-sea
trades studied in most of the maritime optimisation literature
\parencite{psaraftis2013speed}.

In the \CII{} context, \RoRo{} vessels face a distinctive regulatory
environment because the \IMO{} formula uses \GT{} --- a measure of enclosed
vessel volume --- rather than cargo weight or tonne-nautical-miles
\parencite{mepc336, mepc339}.  This means that loading heavier or more cargo
does not improve the \CII{} denominator, and conversely, operating with empty
lanes generates the same \CII{} numerator as a fully loaded voyage (given the
same speed and distance).  Speed and distance are therefore the only levers
available to the operator, making the schedule constraints particularly
consequential.

Ice conditions in the Baltic Sea add further complexity during winter months.
Ice class designations (such as the 1A classification of M/V~Obbola) attract
a correction factor $f_i$ in the \CII{} formula \parencite{mepc336}, partially
offsetting the emissions attributed to ice navigation.  The case study in this
thesis uses the confirmed ice-class factor for M/V~Obbola ($f_i = 1.00$, as
the vessel's ice class provides no correction under current \IMO{} rules for
this vessel category) and does not model seasonal speed variation, focusing
instead on the structural schedule constraints that bind year-round.

% ── 2.5 Research gap ──────────────────────────────────────────
\subsection{Summary and Research Gap}
\label{sec:lit:gap}

The reviewed literature demonstrates growing attention to
carbon-constrained operational planning, yet several gaps remain.
First, existing models predominantly focus on deep-sea or container
shipping; the short-sea \RoRo{} context has received limited attention.
Second, the \GT{}-based \CII{} formula for \RoRo{} vessels is rarely
discussed, and its implication --- that cargo has no direct effect on
\CII{} --- is not reflected in most operational models. Third, the
interaction between fixed port schedules and \CII{} feasibility has not
been systematically analysed. This thesis addresses these gaps through
a model and case study calibrated to Baltic Sea \RoRo{} operations.
